\section{Finite narrowness and scalar projections}
\label{sec:finite-narrowness}
We treat a tangent parent with positive finite narrowness.  Scalar
projection first supplies full-time graphs.  Wide projected intervals
produce a time improvement; the other case forces a horizontal affine
model, whose critical rates are treated in \Cref{sec:critical-rates}.
Throughout this section, \(K=N^{o(1)}\) is a tangent subpower factor
which may increase, and dense means mass at least \(K^{-1}\) in the
stated normalization.  Slowly increasing lists of requirements are
handled only after each fixed finite list has been established.
Conversion to original-resolution outputs is still supplied by
\Cref{prop:candidate-upgrade}.


\begin{remark}[Reference weights and current masses]
\label{a06:summable-pruning}
We use the reference-weight pruning of
\Cref{a03:finite-state-conditioning} in the following scaled form.
Suppose reference weights \(R_e\) satisfy
\(\sum_eR_e\leq K^{(0)} M\), while a current restricted law has
mass at least \(M/K^{(2)}\).  If an appended datum has at most
\(K^{(3)}\) possibilities for each \(e\), deleting entries of
current mass less than \(R_e/K^{(1)}\) loses at most
\[
 \frac{K^{(3)}}{K^{(1)}}\sum_eR_e
 \leq \frac{K^{(0)}K^{(3)}}{K^{(1)}}M .
\]
Choose the subpower \(K^{(1)}\) after \(K^{(0)},K^{(2)},K^{(3)}\) to make this
negligible.  If, before normalization, the numerator at such an
entry and a velocity bin \(b\) is at most \(K\pi(b)R_e\), the
retained entry's conditional velocity probability is at most
\(KK^{(1)}\pi(b)\).  Summing over entries gives the same type of
bound at their common appended state.
These are current floors obtained at a specified pruning stage.
Older witnesses may separately be stored for geometric counts.
Further restriction preserves an unnormalized upper bound, but
does not by itself preserve a normalized conditional bound or a
current lower floor.  A new use of either is justified by the
displayed summable calculation.
\end{remark}


\subsection{A scalar projection theorem at finite meshes}

\begin{theorem}[Scalar mesh projection]\label{thm:scalar-mesh-projection}
Let \(\nu\) be a probability law on trajectories
\((t,y(t),x(t))\), where \(y\) is scalar affine, \(x\) is
scalar quadratic, and all coefficients are bounded by \(K=N^{o(1)}\).
Additional trajectory coordinates are permitted.  Let
\(\mu_{\rm act}\le K\nu\otimes dt\) have mass at least \(1/K\).
Fix \(a=N^{-s}\), \(s>0\), and a dyadic mesh
\(\sigma\sim_{\log,N}N^{-M}\) in \((t,y)\), with fixed \(M>s\).
Suppose \(c>0\) and its reciprocal are polynomially bounded and,
for some \(d<1\),
\begin{itemize}
\item each \((t,y)\)-mesh cell and \(x\)-bin of width
\(p=N^{-r}\), \(0\le r\le s\), have active mass at most
\(K\sigma^2cp^d\);
\item the number of occupied \((\sigma,\sigma,a)\)-cells,
multiplied by \(\sigma^2\), is at most \(Kc^{-1}a^{-d}\).
\end{itemize}
The depth hypotheses use fixed powers, with the rounding and
interpolation conventions of \Cref{sec:framework}.  They imply
\(d\ge0\): otherwise summing the unit-width cap over terminal
support contradicts the dense total mass.

After a dense restriction and passage to a subsequence, there are
labels with disjoint trajectory assignments and a common scale
\(\sigma/K\le\beta\le K\), up to subpower factors, such that
each label has
\begin{itemize}
\item a moving affine interval of width \(K\beta\) containing
its assigned \(y\)-trajectories throughout unit time;
\item a quadratic \(F(t,y)\) with \(|x-F(t,y)|\le Ka\)
throughout unit time on its assigned trajectories;
\item assigned active mass at least \(ca^d\beta/K\).
\end{itemize}
The quadratic norm is at most \(K\) in coordinates adapted to the
moving interval of scale \(\beta\).

If, at \(r=s\), appending a \(y'\)-bin of width \(N^{-u}\)
gives the mass ceiling a further factor \(KN^{-S_1u}\) for
\(0<u\le u_1\), with fixed \(S_1,u_1>0\), then
\(\beta\sim_{\log,N}1\).

\end{theorem}
\begin{proof}
\leavevmode\par
\paragraph{Smoothing the mesh law.}
Let \(\nu\) be the given trajectory prior and
\(\mu_{\rm act}\le K\nu\otimes dt\) the original active incidence
measure. Round the five coefficients of \(y,x\) on a mesh
\(\sigma'\ll\sigma a^2\) by a fixed power. Index the positive-weight
rounded vectors by tags \(j\), and let \(\omega_j\) be the
\(\nu\)-weight of tag \(j\). Mark a dyadic time \(\sigma\)-bin
\(I\) for this tag when
\[
 \mu_{\rm act}(j\times I)\ge \omega_j\sigma/K_1.
\]
Write \(M_j\) for the union of its marked bins. The unmarked
aggregate mass is negligible when the subpower \(K_1\) is large
enough.

Let \(\rho_\sigma\) be the product of five independent uniform
noise laws of width \(\sigma\). The new trajectory index is
\((j,\mathbf u)\), and its coefficients are the rounded vector
of tag \(j\) plus \(\mathbf u\). Define the noisy prior and its
marked incidence measure by
\[
 \nu_{\rm jit}=\sum_j\omega_j\delta_j\otimes\rho_\sigma,
 \qquad
 d\mu_{\rm mark}(j,\mathbf u,t)
   =\mathbf1_{M_j}(t)\,d\nu_{\rm jit}(j,\mathbf u)\,dt.
\]
Thus the noise distribution is the same for every tag; the
independence is between its five coordinates. The domination
\[
 \mu_{\rm act}\{(j,t):t\in M_j\}
 \le K\sum_j\omega_j|M_j|=K\mu_{\rm mark}(\mathrm{all})
\]
shows that \(\mu_{\rm mark}\) has dense mass. This is a tempered
single-world exact problem in base \(N\), version 1 with hidden
\(w=x\); the mark and tag tables have polynomial sizes.

For the marked noisy law, the exact-base density in an \(x_p\)-bin
is at most \(Kcp^d\).  Fix \((t,y)\).  The marked tags that can
contribute have total prior weight at most \(K\sigma cp^d\):
their original incidences satisfy the prescribed position constraints
enlarged by \(K\sigma\), because motion over a time bin and the
coefficient perturbation both cost at most \(K\sigma\).  Summing
the original cell caps and using the mark threshold gives this weight
bound.  The noisy intercept of \(y\) contributes density at most
\(K/\sigma\), proving the exact-base ceiling.

The independent slope noise gives a further factor
\(K\min(1,\rho/\sigma)\) when \(y'\) is confined to an interval
of width \(\rho\).  If the optional angular estimate is assumed,
the same mark calculation at terminal \(x\)-width gives
\[
 Kca^dN^{-S_1u}
\]
for the joint exact-base density with a velocity bin of width
\(N^{-u}\), for every tested \(u<M\) in its range.  The slope
enlargement is much smaller than that bin.  All these unnormalized
ceilings survive further restrictions.

The noisy terminal support, measured by base Lebesgue measure and
counting \(x_a\)-bins, is at most \(Kc^{-1}a^{-d}\).  Every
marked noisy cell is among \(K\) neighbors of a cell carrying
original active incidences.  Homogenize the profiles by the counting
rules of \Cref{sec:framework}, and pass to a subsequence on which
\(C_0=\lim\log_Nc\).  Dense mass, the terminal ceiling, and the
support bound give
\[
 n(s)\sim_{\log,N}ca^d,\qquad f(r)\ge dr-C_0,
 \qquad f(s)-f(r)\le d(s-r).
\]
If the angular estimate was assumed, the resulting problem belongs
to \(C(d)\); choose its positive angular depth range below \(M\).
Then \Cref{a05:scalar-zero-horizon} gives true packets with
arbitrarily small horizon exponents on subsequences.  Without that
estimate, we first obtain it by cropping the slope and intercept.

\paragraph{Obtaining an angular gain.}
When no angular gain is assumed, also homogenize the joint profiles
\(f(r,u)\) with velocity bins through \(U=4(M+s+1)\).  Put
\(R=(1-d)/2>0\), fix a sufficiently small \(\epsilon>0\), and
maximize
\[
 (1-R)r-f(r,u)+\epsilon u,
 \qquad (r,u)\in[0,s]\times[0,U].
\]
At \((s,0)\) its value is \(C_0+Rs\).  The slope-noise ceiling
bounds the objective above by
\(C_0+Rr+\epsilon u-(u-M)_+\).  Continuity of the profiles
therefore gives a maximizer \((r_0,u_0)\) with
\[
 s-r_0\le\epsilon U/R,\qquad
 u_0\le M/(1-\epsilon),\qquad
 f(r_0,u_0)\le f(s,0)+\epsilon U.
\]
In particular \(r_0>0\) and \(u_0+r_0\le U\).  Optimality
with \(u=u_0\) fixed gives a terminal slope cap \(1-R\) ending
at \(r_0\).  Optimality with \(r=r_0\) fixed gives
\[
 f(r_0,u_0+b)-f(r_0,u_0)\ge\epsilon b,
 \qquad 0\le b\le r_0.
\]

Crop the noisy trajectories by the intercept and velocity of \(y\)
at dyadic width \(G\) of order \(N^{-u_0}\).  Make each crop
\(\mathcal C\) a world with its conditional prior and world weight
\(p_{\mathcal C}\), its prior probability.  Subtract the central
affine motion and divide \(y\) by \(G\).  There are polynomially
many crops.  Discard extremely light crops, renormalize their world
weights if necessary, and homogenize
\(p_{\mathcal C}=N^{-\alpha+o(1)}\) on the used worlds.
The resulting data remain tempered and have dense mass: the reciprocal
crop probabilities and the affine coordinate changes have polynomial
bounds.

The new pure density exponent is
\[
 f'(r)=f(r,u_0)+u_0-\alpha.
\]
Indeed, given exact \((t,y)\) and the old velocity bin, the crop
label has only boundedly many possibilities, including for its
intercept since \(y(0)=y(t)-ty'\).  Refinement by that label
preserves the typical old exponent, while the within-world density
is multiplied by \(G/p_{\mathcal C}\).  Summing exceptional
masses with the new world weights costs only a subpower factor.
The old and new velocity grids at corresponding joint precisions
boundedly overlap once the crop is fixed, as in
\Cref{lem:normalization}.  Thus the joint formula at \(r_0\)
and relative velocity depth \(b\) replaces \(f(r,u_0)\) by
\(f(r_0,u_0+b)\).  The cropped problem belongs to \(C(1-R)\)
with angular exponent \(\epsilon\).  Apply
\Cref{a05:scalar-zero-horizon} again to obtain arbitrarily small
horizon exponents.

Lifting such packets to original noisy coordinates, write \(\beta\) for their spatial width scale including the \(G\) factor. Their mass per block time length within the single-world noisy law is at least
\[
 K^{-1}p_{\cal C}\,N^{-f'(r_0)}(\beta/G)
 \ \ge\ \beta c a^d N^{-\epsilon U-o(1)}
\]
by the floor in \Cref{def:true-chart} and the profile formulas. They have an explicit quadratic fit over that time block at thickness \(N^{-r_0+o(1)}\) (the spatial frame change was affine in \(t,y\)). Disjointness over worlds/labels lifts to disjointness within each block using the crop assignments. Total selected success lifts to dense total noisy mass by the choice of world weights (any overall renormalization for discarded worlds here costs at most a subpower).

\paragraph{Full-time fits and return to the original law.}
In the cropped construction, let \(\epsilon\) and the horizon
exponents tend sufficiently slowly to zero.  In the optional angular
case, use the terminal systems obtained without cropping and let only
the horizon exponents decrease.  At each fixed tolerance, first take
samples sufficiently late along the corresponding successful
subsequence.  The resulting diagonal has dense noisy incidence mass
in true assignments with block length \(q\ge1/K\), interval scale
\(\beta\le K\), per-label mass at least
\(qca^d\beta/K\), and quadratic error \(Ka\) throughout the
assigned block.  These are tangent subpower conclusions.

The interval scale satisfies \(\beta\ge\sigma/K\).  The block
bound confines \(y'\) to width \(K\beta/q\le K\beta\)
about the moving-center slope.  At exact base the graph test uses
only \(K\) bins of \(x_a\).  Integrate the exact-base ceiling,
including its additional slope-noise factor, over base area
\(Kq\beta\), and compare with the mass floor.  This forces the
asserted lower bound.  In the optional angular case the same
comparison with the augmented ceiling, at a sufficiently small
positive \(u\), excludes every fixed-power deficiency of \(\beta\)
from one.

The quadratic predictor also has norm at most \(K\) in moving
interval coordinates.  Normalize time to the block and use
\((y-y_c(t))/\beta\) for the spatial variable.  Since
\(|x|\le K\), the mass floor and base ceiling show that
\(|F|\le K\) on a set of Lebesgue measure at least \(1/K\)
in a \(K\)-bounded box.  Polynomial Remez bounds its norm.
As \(q\ge1/K\), extrapolation extends this bound to unit time
at cost \(K\).  Extrapolate the affine position bounds and the
quadratic fit errors along each assigned trajectory in the same way.
The moving center is bounded by \(K\) as well, because an occupied
chart contains a bounded assigned trajectory.

Keep one block's assignments carrying dense selected
\(\mu_{\rm mark}\)-mass; there are only \(K\) blocks since
\(q\ge1/K\). Their floors give
\(\sum_{\rm labels}ca^d\beta\le K\). Averaging over the noise
coordinate \(\mathbf u\) in the displayed product prior fixes
\emph{one common coefficient jitter vector} for which the assigned
marked-time mass, measured by \(\sum_j\omega_j\delta_j\otimes dt\),
is still dense. Assign every original trajectory of a selected tag
to that tag's label in this block.

Undo this common jitter in every label. If its affine and quadratic
components are \(\Delta y(t),\Delta x(t)\), use center
\(y_c-\Delta y\) and predictor
\(F(t,y+\Delta y(t))-\Delta x(t)\). The bounded norms and
full-time fits persist for the original trajectories: their remaining
rounding errors are \(O(\sigma')\), negligible in moving coordinates
compared with \(a\) because \(\beta\ge\sigma/K\).

It remains to return from marked-time mass to \(\mu_{\rm act}\).
Every marked tag-bin that contributes at the fixed jitter has
original active mass at least \(\omega_j\sigma/K_1\). Restore its
entire original active part. Full-time extrapolation permits this
also for bins meeting block edges. The mark threshold therefore
retains at least \(K_1^{-1}\) times the successful marked-time
mass, so the restored original mass is dense. Finally discard
labels whose restored mass is too small compared with
\(ca^d\beta\), using an enlarged subpower denominator and the
summed numerator bound above. This costs negligible mass. Further
dyadic pigeonholing makes \(\beta\) common up to subpower factors
and proves the assertion.

Subpower conclusions over an ensemble of eligible models with common exponent bounds can be taken uniformly: a fixed-power failure along a sequence of individually chosen models would contradict existence with subpower losses on a further subsequence. Equivalently use the worst infimum log-loss over models at each stage then diagonalize. \Cref{thm:scalar-mesh-projection} can thus be used anew on dense restrictions meeting its hypotheses.
\end{proof}
\subsection{Projection of a finite-narrowness parent}\label{a06:projected-parent}
Assume now \(0<\ell<\infty,\ y=(Z,Y)\). Use tangent units and
\[
 v=Z',\quad V_1=(1,v),\qquad a_s=N^{-s},\quad h_s=N^{-ks},\quad
 g_s=N^{-\ell s},\quad b_s=h_s g_s .
\]
The true label \(Q_s\) specifies a block of length
\(q_s\sim_{\log,N}h_s\) and a tilted box
\[
 |Y-Y_Q(t,Z)|\le Kb_s,\qquad |Y'-A_QV_1|\le Kg_s,
\]
where \(Y_Q\) is affine with derivative row \(A_Q\) and
coefficients bounded by \(K\).  By the matrix description in
\Cref{prop:tangent-palette}, unit packet coordinates are
interchangeable, up to \(K\), with translated time and \(Z\)
divided by \(q_s\), and
\(U=(Y-Y_Q(t,Z))/b_s\).  These changes also preserve the stable
local graph comparisons.  Occupancy bounds the speed of the major
center by \(K\), so constant major centering within a block suffices.

Start with any substantial residual and choose a typical parent
with dense success and the prepared upper bounds.  A candidate in
such a parent suffices.  Use its restricted incidence measure,
dominated by \(K\nu\otimes dt\), and retain the parent palette
\(\pi\).  The joint cap for
\((t,Z,Y,X_p,[v]_\chi)\) is \(Kp^d\pi([v]_\chi)\) per
unit base volume at the required finite meshes.

We will integrate these caps over moving triangular-coordinate
boxes whose inverse widths and distortions have fixed-power bounds.
Do so on much finer base meshes, holding the velocity bin fixed.
A scalar-bin offset may vary with the base, provided its variation
in one such cell costs at most \(K\) bin widths; use the prepared
lists for those tests.  Conditional velocity domination by
\(K\pi\) at deep-label and appended-point states is renewed by
deleting states too light relative to their earlier reference
weights.  The lower masses below use the same parent normalization.

\paragraph{The projected groups.}
Choose a sufficiently fine fixed-power \(\chi\), and record
\[
 \alpha=([v]_\chi,[Z(0)]_\chi),\qquad
 Z=Z_\alpha(t)+O(K\chi),\qquad
 w_\alpha=\pi([v]_\chi)\chi.
\]
Here \(Z_\alpha\) uses representative affine coefficients and
\(w_\alpha\) is a weight.  At depth \(s\), also crop
\((Y(0),Y')\) at width \(g_s\) into bins \(\gamma\).
Let \(p_{\alpha\gamma}\) be the prior mass of the joint group.
Condition its prior, divide its active measure by
\(p_{\alpha\gamma}\), and normalize \(Y\) by subtracting the
representative affine motion and dividing by \(g_s\).
We shall apply \Cref{thm:scalar-mesh-projection} with
\[
 c=\frac{w_\alpha g_s}{p_{\alpha\gamma}},\qquad a=a_s.
\]
Use a time and normalized-\(Y\) mesh
\(\sigma\sim N^{-M}\), with \(M>s\) sufficiently large in
terms of \(s,k,\ell\), and take \(\chi\) much finer.
Before conditioning, integration of the joint cap while \(Z\)
tracks its specified motion bounds a projected cell and an
\(X_p\)-bin by
\[
 Kp^d\sigma^2g_s\chi\,\pi([v]_\chi)
   =K\sigma^2w_\alpha g_sp^d.
\]
This is the required projected density bound.

\paragraph{The projected support.}
Fix a true \(Q_s\) and a velocity bin.  Its major width and time
extent allow at most \(Kh_s/\chi\) intercept bins for \(Z(0)\).
For each there are at most \(K\) compatible \(\gamma\)'s:
the row \(A_Q\) determines the slope to \(Kg_s\), and
\(Y_Q(t,Z_\alpha(t))\) determines position to \(Kb_s\) at
an occurrence.  In one group the contributed normalized
\((t,Y)\)-support has mesh-enlarged area at most
\(Kq_sb_s/g_s\).  Its strip slope is bounded by \(K\).

The stable graph lists of \(Q_s\) allow only \(K\) terminal
\(X_{a_s}\)-bins per mesh cell.  Indeed the base uncertainty in
packet axes is at most \(K(\sigma+\chi)/b_s\), which is
negligible even compared with \(a_s\); derivatives are bounded
by \(K\) on each stable ball.  Weighting the support measures,
including these scalar bins, by \(w_\alpha g_s\) and summing
therefore costs at most \(Kq_sh_sb_s\) per true label.  Since
\[
 \nu_Q\sim_{\log,N}a_s^dh_sb_s,
 \qquad \sum_Qq_s\nu_Q\le1,
\]
the total cost is at most \(Ka_s^{-d}\).

Discard groups whose weighted support exceeds
\(p_{\alpha\gamma}a_s^{-d}\) by a sufficiently large subpower
factor; their total prior mass is small by this sum.  Also discard
groups with poor conditional success.  For the remaining nonempty
groups, the support inequality and dense conditional mass with the
unit cap give fixed-power bounds on \(c\) and \(c^{-1}\).
Both hypotheses of the scalar mesh theorem now hold.

We obtain labels \(\Phi_s\) with disjoint trajectory assignments
within each \((\alpha,\gamma)\)-group and across those groups,
carrying unconditioned
active mass at least \(w_\alpha a_s^d\beta_s/K\).  Here
\(\beta_s\) is the interval scale in parent \(Y\)-coordinates,
and throughout unit time
\[
 g_s\sigma/K\le\beta_s\le Kg_s,
 \qquad |X-F_{\Phi_s}(t,Y)|\le Ka_s.
\]
The positions and slopes lie in the corresponding affine packets,
and \(F_{\Phi_s}\) has norm at most \(K\) in their coordinates.
The scalar theorem supplies uniform subpower losses over the eligible
groups.  Homogenize the scale exponents by dense selection on
subsequences.  At finitely many depths, prepare these outputs
successively on dense restrictions, renewing light marginal floors
as needed; the earlier budget
\(\sum_{\Phi_s}w_\alpha a_s^d\beta_s\le K\) keeps the label
counts available.  For \(s\) in a compact positive range, the
exponent \(M\) may be common and does not depend on the finer
choice of \(\chi\).


To compare the two kinds of label, subdivide the true time blocks
into dyadic intervals of length
\[
 q\sim_{\log,N}\min(h_s,\beta_s/g_s),\qquad m\sim_{\log,N}qg_s.
\]
A domain \(D\) specifies one such interval, a \(Y'\)-bin of width
\(g_s\), and a \(Y\)-bin of width \(m\) at its midpoint; its
moving-center slope is the binned value of \(Y'\).

\begin{lemma}[Comparison on common stable domains]\label{a06:projected-comparison}
For the true labels \(Q_s\) and projected labels \(\Phi_s\)
just constructed, use the domains \(D\) above.
After a dense restriction their graph predictions agree to \(K a_s\)
on every retained common comparison ball.  Given \((\alpha,D,Q_s)\)
there are at most \(K\) retained \(\Phi_s\), and conversely at most
\(K\) true labels occur given \((\alpha,D,\Phi_s)\).
List counts use stored marginal witnesses on the common domain;
no floor is asserted for every index in the current pair law.
\end{lemma}
\begin{proof}
We put both lists on common domains before pairing their graphs.  This
allows each graph to keep its own marginal witnesses for the later list
count, even when its mass in the current pair law is small.

\paragraph{Common domains and marginal witnesses.}
Fix \(\alpha\).  For entries in either list use the reference weight
\[
 T_D=q w_\alpha a_s^d m.
\]
These weights have a summable budget.  For a true label \(Q_s\),
sum over its subblocks, the possible \(\alpha\)'s, the \(K\)
slope options per \(\alpha\), and the \(K b_s/m\) position slices.
The result is at most \(Kq_s\nu_Q\).  For a projected label
\(\Phi_s\), there are \(K\beta_s/m\) position slices per
subblock and \(K\) slope choices, giving a total at most
\(K w_\alpha a_s^d\beta_s\).  Summing over either list gives
\(\sum T_D\le K\).  Thus \Cref{a06:summable-pruning} lets us
discard light entries and retain marginal witnesses of mass at least
\(T_D/K\), on any dense restriction.

Evaluate the local graphs for \(Q_s\) at \(Z_\alpha(t)\).
They still predict \(X\) to \(Ka_s\) at the active points.
The affine map from normalized \(D\)-coordinates to the true packet
axes has linear part bounded by \(K\): in the transverse row,
\[
 |Y'_D-A_Q(1,v_\alpha)|\le Kg_s,\qquad qg_s,m\le Kb_s.
\]
Replacing \(Z\) by \(Z_\alpha(t)\) changes packet coordinates
by at most \(K\chi/b_s\), so the evaluation error is negligible
at width \(a_s\).

Choose, before pairing any graphs, a grid in normalized \(D\)
with step much smaller than the common inverse-subpower stable
radii and their fixed holomorphic margins.  The enlargement of a
cell containing an occurrence then lies in that occurrence's stable
neighborhood.  Index each sheet by the cell and its stable
ball/branch choice; there are \(K\) subentries.  All witnesses
for a subentry lie in its common comparison ball.  Delete light
marginal subentries with threshold \(T_D/K_1\), enlarging \(K_1\)
so that the deletion is negligible.  The reference sum remains
subpower, including on a putative dense set of bad pairs.
This construction does not form pairwise intersections of
separately chosen balls and hence does not multiply witnesses by
the number of partners.

On these domains both kinds of graph are algebraic of bounded
degree and have norm at most \(K\).  For \(\Phi_s\), this uses
polynomial extrapolation over transverse variation at most
\(K\beta_s\).  The original joint cap, integrated across the
\(\alpha\)-range of \(Z\) and using the graph test at width
\(Ka_s\), bounds each index-base marginal on sufficiently fine
comparison cells by \(KT_D\) times normalized cell area.  The
Jacobian to physical \((t,Y)\) is at most \(Kqm\).
At still finer meshes these bounds follow from the true-base caps
already prepared above.  Jitter the common paired event point
uniformly within its tiny comparison cell.  Stable margins and
bounded derivatives preserve agreement of values to \(Ka_s\),
and the resulting index-base marginals satisfy the same exact
density ceilings.

\paragraph{Matching with separate reference weights.}
Take norms and scaled center jets on the common comparison balls,
with their fixed holomorphic margins.  Suppose that graph discrepancies
exceed \(a_s\) by a fixed power
on dense pair mass.  Pigeonhole their norm at a dyadic scale
\(M_0\), and group pairs within each \((\alpha,D)\)-comparison
subbox by neighboring center-jet bins at that scale, through a
sufficiently large fixed order.  By \Cref{lem:algebraic-norm},
partner bins are boundedly many neighbors, and all graphs in a
group lie within \(O(M_0)\) of a common graph.  An index enters
boundedly many groups at a fixed discrepancy scale, since it has
only one jet bin.  There are \(O(\log N)\) discrepancy scales and
\(K\) stable subentries, so the reference budget is still at most
\(K\).

Each color has at most \(K(M_0/a_s)^d\) indices in the group.
Indeed, all stored witnesses for one index lie in the comparison
subbox and within \(KM_0\) in \(X\) of the common graph.  Sum
their floors \(T_D/K\) under the cap at that width, using
disjointness within each list, or its subpower multiplicity.
For subpower-large \(M_0\), the unit cap gives the same bound.

Let \(B\) be a group's current bad-pair mass, and let
\(W_c=\sum_{i\in\mathcal I_c}T_i\) be its reference total in
color \(c=1,2\).  Since the total bad mass is dense and the
reference budget is subpower, one group satisfies
\[
 B\ge (W_1+W_2)/K.
\]
Restriction preserves the index-base ceiling \(f_i(z)\le KT_i\).
With the separate color priors \(p_i=T_i/W_c\), normalization of
the pair law gives
\[
 \frac{f_i(z)}{B}
 \le K\frac{W_c}{B}\frac{T_i}{W_c}\le Kp_i.
\]
Integrating the same ceiling also gives \(B\le KW_c\).
Thus the separate reference totals are comparable up to subpower
factors; individual surviving index masses need no lower bound.
Apply \Cref{a02:local-graph-matching} with \(R=M_0/a_s\) to
exclude this group.  Letting the fixed power tolerance decrease
slowly gives discrepancy at most \(Ka_s\) on the retained pairings.

For fixed \((\alpha,D,Q_s)\), every retained projected test is
now within \(Ka_s\) of one of the \(K\) graphs of \(Q_s\) on
the appropriate comparison subbox.  Sum its stored marginal floor
under the cap at width \(Ka_s\).  There are at most \(K\)
possible \(\Phi_s\).  The same argument with the colors reversed
gives the converse list bound.  Summing over stable subentries and
subboxes costs only \(K\); no comparison between weights of
neighboring palette bins is used.
\end{proof}

\subsection{Three wide moving intervals}

\begin{proposition}[Wide-interval improvement]\label{prop:wide-interval}
Use the finite-narrowness parent data of
\Cref{a06:projected-parent}: the joint caps, conditional palette
domination at fine point and path states, and the true \(Q_s\)
graph lists.  The observable \(X\) may be the native signal or
any other bounded quadratic-along-lines signal with these data.
At scales \(s,s+u,s+c_2\), put
\[
 a=a_s,\qquad a_1=aN^{-u},\qquad a_2=aN^{-c_2}.
\]
Suppose that simultaneous projected labels
\(\Phi,\Phi_1,\Phi_2\) of the type constructed above, including
their mass bounds, occur on dense mass and use the same
\(\alpha\)'s.  Let \(\beta_{\min},\beta_{\max}\) be their
smallest and largest interval scales.  Assume
\(0<c_2<(1-d)u/4\) and
\[
 \beta_{\max}\le Kg_s,\qquad
 b_s/\beta_{\min}+h_s\beta_{\max}/\beta_{\min}
       \le N^{-(C_d+3)u},\qquad C_d=10/(1-d),
\]
up to subpower factors, with common scale exponents.  Take all
comparison meshes, including the intercept mesh \(\chi\),
sufficiently fine in fixed powers relative to these widths and
\(h_s,b_s,a_1\), so the transverse evaluation errors are
negligible even at width \(a_1\).

Then some true \(Q=Q_s\) admits a quadratic predictor in the
full base \((t,Z,Y)\) for \(X\), of error \(Ka_2\) throughout
its block on assigned trajectories, carrying original active mass
at least
\[
 q_s\nu_Q(a_2/a)^d/K.
\]
Here \(\nu_Q\) is the full trajectory assignment weight.

\end{proposition}
\begin{proof}
We first show that the finer projected fits differ from a coarse fit
by nearly constant offsets on a true block.  We then find common
nonconstant transverse coefficients, subtract them, and apply scalar
projection in \((t,Z)\).

Use first the short domains in \Cref{a06:projected-comparison}, now at the \(Q_s\) block length and width \(b_s\), slope bins \(g_s\). Given \(Q=Q_s,\alpha\) there are only \(K\) relevant short domains, and only \(K\) possible retained \(\Phi\)'s across this interval. The reverse count given \(\alpha,D,\Phi\) costs \(K\) as well. The \(\Phi_i\) sliced entries have marginal floors \(K^{-1} q_s w_\alpha a_i^d b_s\), by the same width slicing since their slope uncertainties are \(\le K g_s\).

\paragraph{Comparing the projected fits at two domain sizes.}
Compare \(\Phi_i\) and \(\Phi\) on long domains of time length
\(\rho\sim_{\log,N}\beta_{\min}/\beta_{\max}\), slope-bin
width \(\beta_{\max}\), and moving spatial width \(\beta_{\min}\).
Let \(A\) be the common time-times-width Jacobian.  An entry of
color \(i\) has reference weight
\[
 T_i=A w_\alpha a_i^d.
\]
As in \Cref{a06:projected-comparison}, slicing gives a summable
budget of these weights and retained marginal floors \(T_i/K\).
The individual index-base ceilings are \(KT_i\).

The polynomial discrepancies on these long domains are at most
\(KaN^{C_du}\) after negligible deletion.  To prove this, use
neighboring-jet groups at discrepancy scale \(M_0\) and match
width \(a\), as in \Cref{a06:projected-comparison}.  The list
in color \(i\) has at most \(K(M_0/a_i)^d\) members.  Since
\(a_1=aN^{-u}\) is the smallest thickness, whenever
\(M_0/a\ge N^{C_du}\),
\[
 (M_0/a_i)^d\le(M_0/a)^dN^{du}
 \le(M_0/a)^{d+d/C_d}
 \le(M_0/a)^{(1+d)/2}.
\]
The strict inequality \(d+d/C_d<(1+d)/2<1\) leaves room for
subpower losses.  Normalize the current pair law with the separate
color priors \(T_i/W_i\), using its bad mass \(B\) and the
reference totals \(W_i\) exactly as in the earlier proof.
Thus \Cref{a02:local-graph-matching} excludes any fixed-power
excess over \(aN^{C_du}\).  This argument keeps each color's
factor \(a_i^d\) in its own prior.

On the \emph{whole} short domain of a retained pairing this
difference \(F_{\Phi_i}-F_\Phi\) now varies by \(\ll a_i\):
the normalized diameter needed in the long coordinates measured
from the pairing is at most
\(K(h_s/\rho + b_s/\beta_{\min})\), allowing slope error
\(K g_s\) over the short time. Polynomial bounds on fixed
enlargements suffice. Likewise on the whole fiber box
\(Z=Z_\alpha(t)+O(K\chi),\ Y=Y_Q(t,Z)+O(K b_s)\) in this
time interval the same estimate applies. Thus at retained events
\[
 X=F_\Phi(t,Y)+D_i+O(K a_i)
\]
where \(D_i\) can be a constant determined by
\(Q,\alpha,\Phi,\Phi_i\) (evaluate the difference at the block
center on the specified central fiber); \(|D_i|\le K a\).

There are at most \(K(a/a_i)^d\) possible \(\Phi_i\)'s and
hence constants given \(Q,\alpha,\Phi\): on each relevant short
domain the participating \(\Phi_i\) have their tests within
\(K a\) of \(F_\Phi\), so their marginal floors sum by the cap.
Conversely given a short domain, \(\alpha,\Phi_i\), there are
only \(K\) possible paired \(\Phi\).

The reverse bound and the true-label comparison list give only
\(K\) refinements of each sliced \(\Phi_1\) entry. After pruning
light source entries, typical source entries consisting of
\((\alpha,D,\Phi_1,\Phi,Q)\) have active weight at least
\(K^{-1} q_s w_\alpha a_1^d b_s\). Their normalized \((t,Y)\)
density on the short domain is \(\le K\) times uniform by the
finest label test and cap (at sufficiently fine meshes).

Use horizontal coordinates \(z=(t-t_*,Z-Z_*)/q_s\) (constant centering per \(Q\)) and \(U=(Y-Y_Q(t,Z))/b_s\). Expand
\[
 F_\Phi(t,Y)=f_0^\Phi(z)+f_1^\Phi(z)U+f_2^\Phi U^2 .
\]
Here \(f_1^\Phi\) is affine and \(f_2^\Phi\) constant. We seek
one polynomial
\[
 P_{\rm nc}(z,U)=C_1(z)U+C_2U^2,
 \qquad e=aN^{-(1-d)u/2},
\]
with \(C_1\) affine and \(C_2\) constant, such that on dense mass
\((f_1^\Phi(z),f_2^\Phi)=(C_1(z),C_2)+O(Ke)\).
Subtracting it will leave, to error \(Ke\ll a_2\), only the
\(z\)-dependent coefficient and the already counted offsets
\(D_2\). We can then count scalar bins without retaining \(U\)
and project in \((t,Z)\).

The coefficients \(f_1^\Phi,f_2^\Phi\) are bounded by \(K\)
near their active fibers by the moving-interval norm, since
\(Kb_s\) fits inside the interval scale. Evaluation throughout
bounded \(Q\)-coordinates, or subpower enlargements, requires only
fixed-power accuracies: the original coefficients cost at most
\(K\beta_{\min}^{-O(1)}\), and occupied interval centers cost
\(K\). Taking \(\chi\) sufficiently small therefore makes
replacement of \(Z\) by \(Z_\alpha(t)\) negligible.

\par\smallskip\noindent\textbf{Coefficient agreement at a fine state.}\quad

Condition on a fine state consisting of \(Q\) and \((z,U,X)\)
in tiny cells. All comparison precisions, including the direction
grid used below, may be prepared at arbitrarily fine fixed powers.
After light-state deletion, conditional \(v\)-bin weights are
bounded by \(K\pi\). This follows from the \(K\) fine-state
options per end label and does not assert independence from the
projected labels.

Normalize only the whole dense current law to sample a source
event, then sample a target independently conditional on the same
state. Source entries \((\alpha,D,\Phi_1,\Phi,Q)\) retain their
\emph{unnormalized} \(\Phi_1\) ceilings after state deletion;
no floor for their individual surviving masses is needed. Their
reference weights \(q_sw_\alpha a_1^db_s\) sum to at most \(K\).
For the target velocity we instead use the conditional kernel bound
\(K\pi(d_v)\). We claim that the two nonconstant coefficient
pairs agree to \(Ke\) except with negligible sampling probability.

Fix a source entry and a target velocity bin \(d_v\) of width
\(\chi\).  Their joint law is bounded by the source law times
\(K\pi(d_v)\).  We now choose a containing list of target
polynomials before considering their offsets.

First restrict source time to a sufficiently short fixed-power bin.
The fine state ties the source and target base positions much more
closely than \(\chi\), so the target intercept lies within
\(K\chi\) of \(Z_\alpha(t)-tv_{d_v}\).  Throughout this time
bin there are only \(K\) possible target \(\alpha\)'s,
uniformly in \(U,X\).  For each, the whole tilted \(Q\)-box
meets only \(K\) slope/midpoint short domains, and the comparison
list gives \(K\) coarse predictors \(\Phi\) on their union.
A finer target label chooses one of at most \(KN^{du}\) offsets
\(D_1\) for such a predictor.  It does not enlarge the coarse
polynomial list, and its offset does not change the derivative
polynomial.  The source entry already fixes its own polynomial and
offset.

Use source variables \(t,U_\alpha\), \(U_\alpha=(Y-Y_Q(t,Z_\alpha(t)))/b_s\); here \(|U_\alpha|\le K\). Expand \emph{both} candidate polynomials for these variables at source \(Z_\alpha(t)\). Their values with offsets (using \(i=1\) in the offset formula) have to agree within \(K a_1\): fine-state agreement controls \(X,t,Z,Y\) between the events, and in the new expression only \((t,Y)\) are used as arguments of either polynomial. A power-negligible error also covers transport of the two pairs of nonconstant coefficients to this source evaluation from their respective \(z\)'s. All these errors, including those for then freezing the polynomial coefficients on the indicated tiny source time bin, can fit inside \(K a_1\) by the fixed power bounds, \(\chi\) accuracy and still finer bins.

Put \(\tau=(1-d)u/8\).  For each frozen polynomial pair with
nonconstant coefficient gap at least \(e\), the derivative gap
exceeds \(eN^{-\tau}\) outside a set of \(U_\alpha\)-length
\(KN^{-\tau}\).  If the quadratic coefficient difference is
much smaller than \(e/K\), the linear term gives this uniformly;
otherwise apply the sublevel estimate for the affine derivative.
Sum this exceptional length over the \(K\) coarse polynomial
choices.  Offsets introduce no further exceptional sets.

On the complement, monotonicity on a bounded number of intervals
for each polynomial and offset bounds the total length where values
can match to \(Ka_1\) by \(KN^{du}a_1N^\tau/e\).  The two
bounds are therefore
\[
 KN^{-\tau},\qquad
 KN^{du}\frac{a_1N^\tau}{e}
      =KN^{-3(1-d)u/8}.
\]
Both save a fixed power.  The errors in freezing the coefficients
are already within \(Ka_1\), so a coefficient gap of \(Ke\)
at the actual events is covered by these estimates.  Thickening the
at most \(KN^{du}\) intervals by sufficiently fine fixed-power
mesh steps preserves the saving.

Integrate the fine time bins with their lengths.  No factor equal
to their number is incurred.  State deletion only decreases the
source entry's unnormalized ceiling, and normalization of the whole
dense sampling law costs \(K\); there is no division by the
surviving mass of an individual entry.

At such fine source meshes in normalized time and \(U_\alpha\), the source upper weight costs \(\le K q_s w_\alpha a_1^d b_s\) times mesh area per entry, by the joint cap and the \(\Phi_1\) test. This holds with normalization to our overall sampling probability as well. Thus the length tests integrate against this upper source measure and the palette probability with power-small cost times the reference weights; sum those weights. This proves the claimed coefficient agreement.

\par\smallskip\noindent\textbf{A common coefficient field.}\quad
Conditional averaging now assigns a deterministic coefficient pair
to each fine state, agreeing to \(Ke\) with the source pairs on
dense mass. Prune light states to retain velocity domination, then
choose \(Q\) with dense mass relative to \(q_s\nu_Q\), using
\(\sum_Qq_s\nu_Q\le1\). Its normalized \(z\)-marginal is
bounded by \(K\) times area. We next show that these statewise
pairs agree on dense mass with \((C_1(z),C_2)\), for one affine
\(C_1\) and constant \(C_2\), both bounded by \(K\).

Choose first a fixed-power grid step \(\omega\ll\chi\), small
enough that all coefficient evaluation errors below are \(\ll e\).
This choice is independent of the eventual basis, since the
admissible bases and their inverses cost only \(K\). Prepare a
direction grid and the common conditioning-state mesh much finer
than \(\omega\). These direction bins refine the width-\(\chi\)
bins stored in \(\alpha\); each meets only boundedly many of those
older bins.

Now sample from \(\pi\). Each typical state's compatible direction
bins have total palette weight at least \(1/K\); the ball cap gives
separated pairs weight at least \(1/K\). Averaging selects two fine
bins with representatives \((1,v^1),(1,v^2)\), separated to
reciprocal-subpower order, such that dense state mass has a witness
from each. Use these vectors as the axes of the step-\(\omega\)
product grid in \(z\), and assign each state to a nearby vertex.

Along a grid fiber parallel to either axis, its witnesses have only
\(K\) possible \(\alpha\)'s over the entire fiber. Indeed the
fiber prescribes an affine motion of that slope in \((t,Z)\), with
contact errors at most \(Kq_s\omega\). The witness slope bin is
much finer than \(\omega\), so the full-time intercept has only
\(K\) bins at width \(\chi\). The earlier comparison then gives
only \(K\) possible \(\Phi\)'s for this fiber and \(Q\). Their
coefficient pairs restrict to an affine function and a constant on
the fiber. At each compatible state one option on each fiber agrees
with the assigned pair to \(Ke\). Select one polynomial pair per
fiber by averaging; dense state mass satisfies both selected fits.
Keep all matching occurrences, including those with directions
outside the two witness bins.

There are now at least \(K^{-1}\omega^{-2}\) vertices each carrying weight at least \(\omega^2/K\) with that agreement (by the \(z\) upper bound and discarding light vertices), out of a product of side counts \(\le K/\omega\). A dense fraction up to powers of \(K\) of \(3\)-by-\(3\) product grids have all vertices present: the expected proportion of common neighbors of three independently chosen first coordinates is at least the cube of the overall density, by convexity; cube again for three independent opposite coordinates. Each side count is also \(\ge 1/(K\omega)\), so deleting near coincidences (on step-\(\omega\) grids) permits demanding reciprocal-subpower coordinate gaps by taking those cutoffs sufficiently small compared with this density, still with many grids. Fix the first two anchor coordinates on each side so there are \(\ge\omega^{-2}/K\) target pairs completing such grids.

Bilinear interpolation from the four anchors fits the first entry
at the target vertices to \(Ke\); a constant from one anchor fits
the second. Choose the four first-entry values from the selected
affine polynomials on the two anchor rows. The interpolant agrees
identically with those row polynomials. On a target column it is
therefore within \(Ke\) of the selected column polynomial at
both anchor rows. Affine interpolation along that column, using
the inverse-subpower gaps, gives the target fit. The fiber constants
compare through their crossings with one anchor row. Occupied
anchor bounds and the gaps bound both coefficient sets by \(K\).
Grid accuracy transfers the fits to events, and the vertex floors
and target count give dense fitted mass.

It remains to remove the mixed bilinear term. Let \(\mu_Q\) be
the unnormalized law in \(Q\) before the fiber and anchor
selections. At that stage \(Q\) belongs to the fine state with
conditional velocity domination, and \(\mu_Q(\mathrm{all})\le
Km_Q\), where \(m_Q=q_s\nu_Q\). For either selected axis
\(v^j\), this earlier law satisfies
\[
 \mu_Q\{|v-v^1|<r\ \text{or}\ |v-v^2|<r\}
 \leq K m_Q\sum_{j=1}^2\pi(B(v^j,2r))
 \leq K m_Q r^{S'}.
\]
This upper bound survives every subsequent restriction.  If
the final fitted law has mass \(m_Q/K_{\rm ret}\), choose the
reciprocal-subpower radius \(r\) slowly enough that
\(K K_{\rm ret}r^{S'}=o(1)\), first proving the assertion at
fixed power radii.  A dense non-axis residual remains.
The law is dominated by \(K\) times the original trajectory prior
restricted to the full assignment of \(Q\), times Lebesgue measure
in time.  That restricted prior has total mass \(\nu_Q\), so
some indexed line has retained duration at least \(q_s/K\).
Its fixed \(\alpha\) and \(Q\) permit only \(K\) coarse
predictors, so one \(\Phi\) holds for duration \(q_s/K\).
On those retained times the bilinear interpolant agrees to
\(Ke\) with the affine function \(f_1^\Phi\).
In the chosen basis the components of \((1,v)\) are
\[
 \frac{v^2-v}{v^2-v^1},\qquad
 \frac{v-v^1}{v^2-v^1}.
\]
Both are bounded below in absolute value by a reciprocal
subpower.  The quadratic coefficient of the bilinear
interpolant on this line is its mixed coefficient times their
product.  Remez on the retained time set therefore bounds
the mixed coefficient by \(Ke\). Removing that term gives the
claimed affine \(C_1\), while the fitted second entry gives the
constant \(C_2\). The argument used the earlier unnormalized
direction bound throughout.


\par\smallskip\noindent\textbf{Scalar support after subtraction.}\quad
The common field just proved supplies the proposed
\(P_{\rm nc}=C_1(z)U+C_2U^2\). Subtract it from \(X\) in \(Q\).
The resulting quadratic signal is subpower bounded on assigned
trajectories. On a sufficiently fine \(z\)-mesh, with \(U\) no
longer recorded, each cell has at most \(K(a/a_2)^d\) residual
bins of width \(a_2\) after pruning. Indeed, at fixed
\([v]_\chi\), the values are
\(f_0^\Phi(z)+D_2+O(Ka_2)\): there are \(K\) choices for
\(\alpha,\Phi\) and \(K(a/a_2)^d\) for \(D_2\). The error
\(Ke\) fits within this width because \(c_2<(1-d)u/4\).

To sum this list count over velocities, use the output state
consisting of the \(z\)-cell and residual scalar bin in the chosen
\(Q\). It has only \(K\) options per sufficiently deep end label,
since \(P_{\rm nc}\) has fixed-power coefficient bounds. Delete
light appended entries against the earlier end-history weights,
whose sum inside \(Q\) is at most \(Kq_s\nu_Q\). The retained
output states then have conditional velocity domination, hence
compatible velocity bins of palette mass at least \(1/K\).
Integrating the preceding count against \(\pi\) proves the
claimed residual-bin count. The fine \(z\)-mesh determines the
intercept to width \(\chi\) up to \(K\) choices and freezes
\(f_0^\Phi(z)\) to \(Ka_2\); all these meshes are prepared
simultaneously.

Normalize the current mass by \(q_s\nu_Q\).  In unit
\(Q\)-time and scaled \(Z\), the mass per unit projected base
volume in a scalar bin of width \(a_2\le p\le1\) is at most
\(K(p/a)^d\), also with the angular factor from \(\pi\).
For a \(z\)-cell of side \(\sigma_*\), integrating over
\(|U|\le K\) uses old base volume at most
\(Kq_s^2b_s\sigma_*^2\).  On still finer meshes the
\(X\)-bin can follow the offset \(P_{\rm nc}\).
Now use \(\nu_Q\sim_{\log,N}a^dh_sb_s\) and the palette ball
bound at sufficiently small positive exponents.  The scaled velocity
remains \(v\).

The prior for the next projection is the true assignment of \(Q\).
Restrict it, if necessary, to trajectories with bounded coefficients
on the normalized full block.  Every trajectory in the final dense
events has these bounds by construction.  Since \(Q\) was selected
heavy relative to its full trajectory weight, this conditioning costs
only \(K\).  Apply \Cref{thm:scalar-mesh-projection} to
\((X-P_{\rm nc},Z)\), with \(c=a^{-d}\), terminal width
\(a_2\), and its positive angular cap.  The terminal density is
then \(ca_2^d=(a_2/a)^d\).  The support
count just proved is its second hypothesis.  It gives labels of full
horizontal width up to \(K\), hence a quadratic prediction in
\((t,Z)\) throughout the block, with normalized mass at least
\((a_2/a)^d/K\).  Restoring \(P_{\rm nc}\) gives the asserted
quadratic in the full base and the required mass.

These last comparison meshes may have arbitrarily fine fixed-power
widths.  The equations with errors \(O(a_i)\) are evaluated on
true base variables; the prepared pure and joint caps and palette
domination apply to the finer native point states.  No improvement
of those equation errors is needed.  If this construction is used
inside an earlier true parent after affine changes, select heavy
descendants using their full assignments nested in that parent's
prior, whose weights sum to at most that prior in each block.

\end{proof}
For the original parent observable \(X\) this is a time-improvement candidate (subtract the quadratic from its native polynomial). Fix for example \(s=1\); choose \(u>0\) sufficiently small depending on \(s,k,\ell,d\), then \(0<c_2<(1-d)u/4\). Prepare the three projections as above (including the mutual comparison lists with their corresponding true packets where needed). If at all three scales the widths satisfy \(\beta_{s'}\sim_{\log,N}g_{s'}\), \Cref{prop:wide-interval} applies. Otherwise passing to subsequences some (henceforth denoted \(s\)) has
\[
 \beta_s=N^{-B_0+o(1)},\qquad \ell s<B_0\le\ell s+M.
\]
We now deal with this remaining case. In choosing \(\chi\) it suffices to anticipate the fixed upper bounds here and the comparison errors through these scales.

\subsection{The horizontal model}

\begin{proposition}[Reduction to the horizontal model]\label{prop:horizontal-model}
Suppose \(\beta_s=N^{-B_0+o(1)}\) with
\(\ell s<B_0\leq\ell s+M\).  On every substantial residual either
there is a strip-improvement candidate, or a true-parent normalization
gives \(k=\ell\) and \Cref{a06:horizontal-equation}, with fixed
\(e_0>0\).  At sufficiently small prescribed relative depths the
whole true boxes have horizontal scale \(H\) and transverse scale
\(H^2\) in \Cref{a06:centered-horizontal-box}.
Mass is normalized by the chosen parent's full assignment times
its time length.
\end{proposition}
\begin{proof}
\leavevmode\par
\paragraph{From projected intervals to a row list.}
Choose \(r>s\) with \((k+\ell)r>B_0>\ell r\).
Given \(Q_r\), its \(Q_s\)-history, and sufficiently fine
horizontal position and velocity bins, there are only \(K\)
possible \(\beta_s\)-bins for \(Y'\).  Indeed \(\alpha\)
is known up to neighbors.  The \(Q_s,\Phi_s\) comparison domains
then have \(K\) possibilities: \(Q_r\) determines \(Y\) to
\(Kb_r\), while \(Q_s\) determines the \(g_s\)-slope bin.
Translating from the occurrence time to a comparison-subblock
midpoint costs \(Kg_sq\), its prescribed width, and
\(b_r\lesssim_{\log,N}\min(b_s,\beta_s)\).
Thus only \(K\) midpoint bins, and hence \(K\) projected
labels, are possible.  Each label's full-time moving interval
confines \(Y'\) to width \(K\beta_s\).  Take the horizontal
and velocity meshes sufficiently fine also relative to \(\chi\).

Normalize the true parent \(Q=Q_r\).  Translate time and \(Z\)
by constants and divide by \(q_r\), keeping horizontal velocity
\(V_1=(1,v)\).  Subtract \(Y_Q\) and divide the transverse
coordinate by \(q_rg_r\).  The normalized transverse velocity is
\[
 \frac{Y'-A_QV_1}{g_r},
\]
bounded by \(K\), and the preceding list has accuracy
\(E=N^{-(B_0-\ell r)}\).  A sufficiently deep true label has row
\[
 T=(A_{\rm deep}-A_Q)/g_r
\]
whose evaluation \(TV_1\) predicts this velocity to \(\ll E\).
Fix base and velocity meshes \(\xi\ll E\), fine enough for
these lists.

We verify the remaining row-fit hypotheses.  At a state consisting
of the deep row label and fine horizontal location, use sufficiently
deep histories, append this state, and delete light states to obtain conditional
velocity domination by \(K\pi\).  Its supporting velocity bins
have palette mass at least \(1/K\).  The angular cap therefore
supplies pairs separated by a reciprocal subpower, with palette
product weight at least \(1/K\).  The two evaluations of the fixed
row at actual witnesses bound \(T\) by \(K\).

For any fixed such pair at a horizontal cell, the two lists for the
normalized transverse velocity constrain compatible \(E\)-row
bins to \(K\) possibilities by interpolation.  Bounded rows and
velocities make the binned-velocity error negligible.  Choose one
state per row-bin option and sum with the palette product law.
Since every option has witness-pair weight at least \(1/K\), there
are at most \(K\) row options per horizontal cell.

The horizontal joint cell cap, before dividing mass by
\(q_r\nu_Q\), is
\[
 K\xi^2q_r\nu_Q\,\pi(v_\xi).
\]
It follows by summing the joint state bounds with their earlier
weights on which the \(Q\)-base cap holds.  The palette interval
bounds through \(\xi\) follow from the power-scale caps and
interpolation.  Choose a \(Q_r\) heavy relative to
\(q_r\nu_Q\), and use its full conditioned trajectory prior.
All hypotheses of \Cref{a06:row-fit} now hold in the normalized
chart.

\paragraph{The horizontal equation and its scales.}
The row fit gives normalized transverse velocity
\((B+\lambda Jz)\cdot V_1+O(KE)\) on dense mass.
If \(|\lambda|\) is power-small on a subsequence, this is a
strip candidate: multiply by \(g_r\) and restore \(A_QV_1\).
Otherwise \(|\lambda|\sim_{\log,N}1\).  Subtract \(B\cdot z\)
from the transverse coordinate and divide by \(\lambda\).
Reusing \((t,Z,Y)\) for these coordinates, we obtain
\begin{equation}\label{a06:horizontal-equation}
 |Y'-Jz\cdot V_1|\le N^{-e_0}
\end{equation}
for some fixed \(e_0>0\), decreased with room if needed.
The expression on the left before absolute values is constant
along each trajectory.

At a small relative depth \(s'>0\), a child \(Q_{r+s'}\) has
block length \(H=q_{r+s'}/q_r\sim_{\log,N}h_{s'}\),
horizontal diameter at most \(KH\), and a fixed row predicting
transverse velocity to \(KN^{-\ell s'}\).  Compare this row
with \(Jz\) using separated velocities in a typical fine-location
state.  Interpolation gives row error at most
\(K(N^{-\ell s'}+N^{-e_0})\).  These comparisons use
unnormalized domination and earlier weights summing inside
\(Q_r\) to at most \(Kq_r\nu_Q\); delete light refinements
in that normalization.

Choose \(s'\) small enough that the \(N^{-e_0}\) error is
negligible at both relative scales.  After summable deletion, each
used child carries normalized mass at least
\(H\nu_{Q_{r+s'}}/(K\nu_Q)\).  Prior-times-time domination
then forces its active time extent to be at least \(H/K\).
But its fixed row predicts the second component of \(Jz\),
namely \(t\), to \(KN^{-\ell s'}\).  Therefore
\[
 N^{-ks'}\le KN^{-\ell s'},
\]
so \(k\ge\ell\).  If \(k>\ell\), replace \(Jz\) in
\Cref{a06:horizontal-equation} by its value at an active child
center.  The error from the horizontal diameter is smaller by a
power than the nominal velocity width, giving a strip candidate
at \(Q_{r+s'}\).  Restoring earlier coordinates multiplies
velocity error by \(|\lambda|g_r\); dividing the displayed
child mass floor by \(H\) and restoring the \(\nu_Q\)
denominator gives the required candidate mass.  Thus it remains
only to consider \(k=\ell\).

\paragraph{Whole horizontal boxes.}
Choose an occupied center \((z_c,Y_c)\) of a used child and set
\begin{equation}\label{a06:centered-horizontal-box}
 \widehat z=(z-z_c)/H,\qquad
 \widehat Y=(Y-Y_c-Jz_c\cdot(z-z_c))/H^2.
\end{equation}
The true strip row is within \(KH\) of \(Jz_c\), by the
preceding two-velocity comparison at a point within horizontal
distance \(KH\) of the center.  With \(k=\ell\), its true
transverse spatial width is \(KH^2\) up to subpower factors.
The coordinate changes used to obtain the horizontal equation
alter these comparisons only by \(K\).  A slope discrepancy
of \(KH\) over horizontal distance \(KH\) costs \(KH^2\),
so these centered coordinates describe the entire true tilted
box, up to \(K\), rather than only its retained incidences.

The two-velocity comparison can use common-location meshes much
finer than the box widths: the row is fixed, and the variation of
\(Jz\) between the nearby true-velocity witnesses is harmless.
Prepare finitely many requested depths simultaneously, then let
the grids grow slowly, discarding labels and states too light
relative to their earlier weights.  Time is unit time in \(Q_r\),
and the prior is its full true line assignment conditioned to
probability.  Descendant mass floors and deletions continue to use
the full nested assignment weights even if only some assigned lines
participate in the later restrictions.
\end{proof}

\subsection{Affine approximation and residual normalization}

The horizontal equation first forces the local hidden sheets to be
affine to their fitted accuracy.  We then subtract one parent affine
fit.  This gives a bounded residual with caps and state lists at
arbitrarily fine prescribed meshes; those preparations will not require
the horizontal equation at the same accuracy.  Finally, comparison of
nested affine fits gives the slope bounds used in the next section.

\begin{lemma}[Affine approximation of the hidden sheets]
\label{a06:horizontal-affine-approximation}
In \Cref{prop:horizontal-model}, for finitely many sufficiently small
relative depths \(s'\), there are \(K\) affine entries per true label
fitting the hidden signal to width \(Ka_ra_{s'}\) throughout the
corresponding block, including \(s'=0\).  No bound on a sheet's affine
part in the initial tangent coordinates is required.
\end{lemma}
\begin{proof}
Fix one of the required depths, including depth zero.  Until we return
to the parent at the end of the proof, \((t,Z,Y)\) denotes the current
label's centered coordinates from \Cref{a06:centered-horizontal-box},
and \(z=(t,Z)\).  Keep the hidden variable \(w\) in the initial tangent
units of derivative scale one.  Its fitted width is
\(\Delta=a_ra_{s'}\), up to \(K\).  In this centered frame the actual
event direction is
\[
 {\bf W}=(1,v,Jz\cdot V_1+O(KN^{-e_0}/h_{s'})).
\]
Indeed subtracting the central plane and rescaling changes the
horizontal row to \(J(z_{\rm old}-z_c)/H\).  For \(Q_r\) itself the
unit horizontal coordinates of \Cref{a06:horizontal-equation} suffice.
Choose the required depths small enough that the displayed velocity
error is at most \(N^{-e}\), with fixed \(e>0\).

Use the stable-ball preparation of \Cref{prop:tangent-palette} for
the hidden test sheet \(f\), choosing the branch on which the observed
\(w\) lies within \(K\Delta\).  Above the high-curvature cutoff, the
affine quadratic center already gives the required fit.  Below that
cutoff, the discriminant bounds and simple-root stability give balls
with fixed holomorphic margins and inverse-subpower radii.  Passing
from packet axes to the centered boxes costs only \(K\), and the
ball and branch choices number at most \(K\) per true label.

In each ball/branch entry, remove trajectories whose associated time
set has relative length below a sufficiently small reciprocal
subpower.  Charge this deletion to the full true trajectory weights
using prior-times-time domination.  Also discard entries too light
relative to block length times their true assignment weight.  These
thresholds sum within \(Q_r\), so the deletions are negligible.
The true-label base cap then gives inverse-subpower Lebesgue filling
in every retained entry's centered box.  Store the long ball/branch
witnesses for subsequent restrictions.  After the fine-state prunings
below, repeat the light-entry deletion when current base filling is
needed; the earlier long witnesses still give the chord estimates.

Let \(M_f\) be the best sup-norm affine approximation error on the real
comparison ball.  By \Cref{lem:algebraic-norm}, applied with margin
after subtracting an affine function, all derivatives of order at
least two through any fixed required order are bounded by \(KM_f\).
Along a long witnessing chord, the Hessian evaluated twice on the
actual direction \({\bf W}\) is at most \(K\Delta\) at its events.
To see this, take a complex parameter disk around the event of radius
comparable to the ball radius divided by the direction length, with
the fixed margin supplied above.  It contains the earlier
same-ball/branch times of relative parameter measure at least
\(1/K\).  On those times, the sheet differs from the affine line
signal by \(K\Delta\).  The one-variable form of
\Cref{lem:algebraic-norm} bounds its second derivative at the event;
the parameter rescalings cost only \(K\).  This argument imposes no
power bound on \(w\) or \(M_f\).

Use fine states consisting of the path label, fine base position,
and ball/branch choice.  They have \(K\) ambiguity given sufficiently
deep labels.  After the usual light-state deletion, conditional
palette domination supplies three actual velocity witnesses in one
state, separated by reciprocal-subpower distances and carrying the
long-branch tests.  This conditioning uses no uncontrolled
\(w\)-mesh.  Transfer their Hessian estimates to a common base point
and to the horizontal directions \((1,v,Jz\cdot V_1)\).  The velocity
error and sufficiently fine base meshes give an error
\(KM_fN^{-e}\), because only second and higher derivatives enter.

Put
\[
 {\cal X}=\partial_t-Z\partial_Y,\qquad
 {\cal Z}=\partial_Z+t\partial_Y.
\]
For fixed \(v\), the square \((\mathcal X+v\mathcal Z)^2f\) is exactly
the frozen horizontal second directional derivative: the coefficient
of \(\partial_Y\) is \(-Z+tv\), and
\[
 (\mathcal X+v\mathcal Z)(-Z+tv)=-v+v=0.
\]
In particular an affine function contributes zero before the actual
direction is perturbed.  The error above depends on \(M_f\), rather
than on the size of the discarded affine part.

Interpolate the quadratic in \(v\) at the three separated witnesses.
On the base projection of retained events in a well-filled entry,
\[
 \mathcal X^2f,\quad \mathcal Z^2f,\quad
 (\mathcal X\mathcal Z+\mathcal Z\mathcal X)f
 =O\bigl(K(\Delta+M_fN^{-e})\bigr).
\]
These expressions are algebraic analytic functions of bounded
relation degree on the stable enlargement.  The Remez and derivative
bounds in \Cref{lem:algebraic-norm} extend the estimates to the real
comparison ball and to every fixed further derivative needed below.

The central commutator \(\mathcal T=[\mathcal X,\mathcal Z]
=2\partial_Y\) recovers the remaining second derivatives.  Since
\(\mathcal T\) is central,
\[
 \mathcal X^2\mathcal Z
   =\mathcal Z\mathcal X^2+2\mathcal T\mathcal X,
 \qquad
 \mathcal X\mathcal Z\mathcal X
   =\mathcal Z\mathcal X^2+\mathcal T\mathcal X.
\]
Consequently
\begin{align*}
 \mathcal X(\mathcal X\mathcal Z+\mathcal Z\mathcal X)
       -2\mathcal Z\mathcal X^2&=3\mathcal T\mathcal X,\\
 \mathcal Z(\mathcal X\mathcal Z+\mathcal Z\mathcal X)
       -2\mathcal X\mathcal Z^2&=-3\mathcal T\mathcal Z.
\end{align*}
The preceding differentiated bounds therefore control
\(\mathcal T\mathcal Xf\) and \(\mathcal T\mathcal Zf\), and then
\[
 \mathcal T^2f
   =\mathcal X\mathcal T\mathcal Zf-
      \mathcal Z\mathcal T\mathcal Xf.
\]
With \(E_f=K(\Delta+M_fN^{-e})\), all the following expressions are
bounded by \(E_f\), after increasing \(K\):
\begin{align*}
 \mathcal T^2f&=4f_{YY},\\
 \mathcal T\mathcal Xf&=2(f_{tY}-Zf_{YY}),&
 \mathcal T\mathcal Zf&=2(f_{ZY}+tf_{YY}),\\
 \mathcal X^2f&=f_{tt}-2Zf_{tY}+Z^2f_{YY},&
 \mathcal Z^2f&=f_{ZZ}+2tf_{ZY}+t^2f_{YY},\\
 \tfrac12(\mathcal X\mathcal Z+\mathcal Z\mathcal X)f
   &=f_{tZ}+tf_{tY}-Zf_{ZY}-tZf_{YY}.
\end{align*}
Since the comparison coordinates are bounded by \(K\), these
identities bound the entire ordinary Hessian by
\(K(\Delta+M_fN^{-e})\).  Taylor's affine polynomial on the ball
gives
\[
 M_f\le K\Delta+KM_fN^{-e}.
\]
Absorb the second term for sufficiently large \(N\).  No fixed-power
bound on \(M_f\) is needed for this absorption.

The resulting affine fit agrees with \(w\) to \(K\Delta\) on the
stored long ball/branch times of each contributing line.
Their difference is affine on that line, so extrapolation gives
the same bound, up to \(K\), throughout the corresponding block.
There are at most \(K\) fits per true label.  Make these
preparations for each fixed finite family of depths, then let the
family grow slowly, keeping the earlier entry and long-line
witnesses when restricting further.
\end{proof}

\begin{lemma}[Residual normalization and fine-state bounds]
\label{lem:horizontal-affine}
Choose one successful parent affine fit \(L\) from
\Cref{a06:horizontal-affine-approximation} and set
\(x=(w-L)/a_r\).  On the resulting dense restriction, \(x\) is
affine along trajectories and \(|x|+|x'|\le K\).

Let \(\mu\) be the restricted incidence measure in normalized
\(Q_r\) coordinates, with mass divided by \(q_r\nu_{Q_r}\), without
further normalization by its retained mass.  For any prescribed
finite family of fixed-power base, scalar and velocity meshes,
the preparations give
\begin{align}
 \mu\{(t,Z,Y)\in C,\ x\in J_p\}
     &\le K|C|p^d,\notag\\
 \mu\{(t,Z,Y)\in C,\ x\in J_p,\ v\in b_\chi\}
     &\le K|C|p^d\pi(b_\chi).
 \label{a06:residual-caps}
\end{align}
Here \(C\) is a cell in the normalized three-dimensional base,
\(|C|\) is its volume, \(J_p\) is a scalar bin of fixed-power
width \(0<p\le1\), and \(b_\chi\) is a velocity bin at the
prescribed fixed-power width \(\chi\).  The palette \(\pi\) is
the earlier parent palette.  The bounds persist under restriction;
\(K\) may depend on the finite mesh preparation, with increasing
families handled by slow diagonalization.

Sufficiently deep true labels determine these base and residual
bins with \(K\) choices, also beyond the accuracy range of the
approximate horizontal equation.
\end{lemma}
\begin{proof}
Return to the \(Q_r\) coordinates of
\Cref{a06:horizontal-equation}.  The full-block parent fit makes
\(x=(w-L(z,Y))/a_r\) affine and bounded along the used lines, and
therefore also gives \(|x'|\le K\).  We first recover fine residual
states.  Only afterwards will we transfer the velocity numerator
to those states.

Freeze the dense law after choosing \(L\), on which the actual
residual \(x\) is bounded.  At a sufficiently deep true level
\(R\), prune light ball/branch entries against block length times
their full true assignment weights.  In the initial tangent mass
units these references sum within \(Q_r\) to at most
\(Kq_r\nu_{Q_r}\).  The deletion is therefore affordable without
any fine \(x\)-state.  The base ceiling gives inverse-subpower
filling in each retained branch entry's own true axes.  On those
filled base points its analytic residual
\[
 g_R=(f_R-L)/a_r
\]
is bounded by \(K\).  By \Cref{lem:algebraic-norm}, its norm and
fixed derivatives are bounded by \(K\) on the stable comparison
interior, even if \(L\) has enormous coefficients.

Choose \(R\) so that its graph error is below the desired residual
mesh.  A still deeper label supplies the relative location accuracy
required by the derivative bounds.  It then determines at most
\(K\) residual and base bins.  This uses true packet axes and
remains valid beyond the accuracy range of the approximate
horizontal equation; it uses no conditional floor for the new bins.

Now append the recovered state to each earlier end history \(e\).
In initial tangent mass units its reference weight \(R_e\) obeys
\(\sum_eR_e\le Kq_r\nu_{Q_r}\), and the retained velocity
numerator is at most \(K\pi(b)R_e\).  There are at most \(K\)
appended choices per history.  Delete current appended entries
of mass below \(R_e/K_1\), choosing \(K_1\) after the inverse
success fractions for \(Q_r\) and \(L\).  The reference sum
pays for this deletion by \Cref{a06:summable-pruning}.
Division on a retained entry gives conditional velocity domination;
summing with its pre-deletion weight gives the corresponding joint
bound at the appended state.

The pure residual cap is \(Kp^d\) per unit normalized base volume.
Indeed integrate the hidden \(w\)-cap at width \(pa_r\) about the
exact offset \(L\), with the \(Q_r\) Jacobian and prior
denominator.  This bound also holds for the current pre-deletion
weights.  Combining it with the transferred numerator gives
\(Kp^d\pi(b_\chi)\) at the required meshes, proving
\Cref{a06:residual-caps}.  No slow variation of \(L\) on the old
tangent meshes is required.

If a probability prior on bounded participating trajectories is
needed, conditioning away the other trajectories costs at most
\(K\).  The full true descendant assignment weights relative to
\(Q_r\) remain the reference weights for counts and deletions.
In particular, the child target mass per unit block time is
\(a_{s'}^dh_{s'}^3\), up to \(K\).
\end{proof}

For subsequent counts let \(\nu^\flat\) be the full trajectory
prior conditioned on the true assignment of \(Q_r\).  It is not
conditioned on the success of \(L\).  The current measure \(\mu\)
is dominated by \(K\nu^\flat\otimes dt\).  Write \(Q_s\) for
a true child at depth \(s\) relative to \(Q_r\), and
\(H_s=q_{r+s}/q_r\sim_{\log,N}h_s\) for its block length.
The \(H/H^2\) geometry and \(k=\ell\) give
\begin{equation}\label{a06:horizontal-budgets}
 \nu^\flat(Q_s)\sim_{\log,N}a_s^dh_s^3,\qquad
 m_{Q_s}=H_s\nu^\flat(Q_s),\qquad
 \sum_{Q_s}m_{Q_s}\le1.
\end{equation}
These full assignment weights remain the reference budgets after
restrictions of the current incidence law.

\begin{lemma}[Slopes of nested affine fits]
\label{a06:nested-affine-slopes}
Use the residual \(x\) from \Cref{lem:horizontal-affine}, and let
\(s\) denote depth relative to \(Q_r\).  At the required small
depths, the affine predictions \(L_s(t,Z,Y)\) for \(x\) may be
prepared on a dense restriction so that, on active boxes,
\begin{equation}\label{a06:affine-slope-bounds}
 |\nabla_z L_s+(\partial_Y L_s)Jz|
       \le K(1+a_s/h_s),\qquad
 |\partial_Y L_s|\le K(1+a_s/h_s^2).
\end{equation}
The predictions have error \(Ka_s\) throughout their assigned
blocks, with \(K\) subentries per true label.  The coordinates
\((t,Z,Y)\) in this statement are those of the normalized parent
\(Q_r\).
\end{lemma}
\begin{proof}
The affine predictions follow by subtracting \(L\) and dividing
the fits of \Cref{a06:horizontal-affine-approximation} by \(a_r\).
We compare them along a nested depth grid and keep full-domain
witnesses for each pair.  For consecutive depths \(s_{i-1},s_i\),
the paired ancestor and child affine subentries have only \(K\)
options per true child label, whose history and time block are
already specified.  Discard pairs whose mass is below a
sufficiently small subpower fraction of child block length times
the child's full trajectory weight.  These thresholds are
summable.  More explicitly, if the number of allowed pairs is
\(K_{\rm pair}\), deletion below
\(\epsilon m_{\rm child}/K_{\rm pair}\) loses at most
\[
 \epsilon\sum_{\rm child}m_{\rm child}\le\epsilon.
\]
Choose \(\epsilon\) negligible relative to the current dense
mass.  No independent choices of time blocks enter this pair
count.  Each fixed depth grid is pruned first, and the grid is
then allowed to grow slowly.

The child base cap gives inverse-subpower filling for every
retained pair's full-domain witnesses.  Both affine fits are close
to \(x\) there, so affine Remez bounds their difference on the
whole child box by \(Ka_{s_{i-1}}\).  Its slope along the child's
centered horizontal plane is at most
\(Ka_{s_{i-1}}/h_{s_i}\), and its transverse slope is at most
\(Ka_{s_{i-1}}/h_{s_i}^2\).  The same bounds apply at an active
descendant location: that location lies in the enlarged ancestor
box, and \(Jz\) varies by at most \(Kh_{s_i}\) within the
child box.

Start at relative depth zero with the zero prediction and the
bound \(|x|\le K\).  Telescope over a grid with depth gaps
tending sufficiently slowly to zero.  For \(p=1,2\) and
\(s_i\le s\),
\[
 a_{s_{i-1}}/h_{s_i}^{p}
     \le K(1+a_s/h_s^{p}).
\]
The sum of the slope increments therefore gives
\Cref{a06:affine-slope-bounds}.  All losses remain subpower on a
slowly growing grid.  The same witness deletion also prepares
nonconsecutive ancestor/child pairs when needed.
\end{proof}

For later counts, one may replace \(Y\) by
\(Y^0=q_0-tp\), where
\[
 Z=p+tv,\qquad P=(p,q_0)=(Z(0),Y(0)),
\]
with error at most \(N^{-e_0}\), while keeping \(x\) unchanged.
The bounds transfer by bounded enlargements at meshes and
evaluation accuracies safely coarser than that error, including
all slope amplifications.  Conversely, a fit with controlled
coefficients can then use the true \(Y\) in a candidate.
A time improvement in these units lifts by restoring
\(w=L+a_rx\), with hidden derivative one.

\subsection{Full-time coefficient counts in the horizontal model}

For the critical and subcritical rates, the next lemma supplies one
last consequence of the prepared model: when \(k\le1/2\), small
scalar spacetime support concentrates the full-time affine coefficients.


\begin{lemma}[Horizontal coefficient count]\label{a06:horizontal-coefficient-count}
Suppose \(k=\ell\leq1/2\), with the bounded residual
\(x\) of \Cref{lem:horizontal-affine}.  Fix sufficiently small
\(s>0\) and \(D=a_s^u\), with bounded \(u\geq1\).
Group by a sufficiently fine velocity bin \(d_0\) and a
\(D\)-bin of \(P=(Z(0),Y(0))\), and put \(W_G=\pi(d_0)D^2\).
On a dense restriction, typical groups have active and prior mass
comparable to \(W_G\) up to \(K\), and their \((t,x)\)-support
requires at most \(K a_s^{-1-d}\) squares of side \(a_s\).
Their full-time affine coefficient bins can be restricted to
at most \(K a_s^{-d}\) bins of width \(a_s\), each of active
mass at least \(W_Ga_s^d/K\).
After subdividing time at a fixed-power length \(\eta\ll a_s\)
and summably pruning, the group, time subbin and \(x_{a_s}\)
determine a list of at most \(K\) retained coefficient bins.
\end{lemma}
\begin{proof}
Work with the full prior \(\nu^\flat\), current measure
\(\mu\), and true-child budgets in \Cref{a06:horizontal-budgets}.
Restrictions may be taken anew on dense mass in this normalization.
The depths are sufficiently small for the horizontal equation,
centered boxes, and affine preparations, simultaneously at the
finitely many required scales.  When \(k\le1/2\),
\Cref{a06:affine-slope-bounds} and the bounded parent positions
bound every ordinary slope of \(L_s\) by \(K\).  We retain
the joint finite-mesh caps for \((t,Z,Y,x,v)\).

Fix small \(s>0\), \(D=a_s\) (or \(D=a_s^u,\ u\ge1\) bounded, depths small enough), and group by a very fine \(v\)-bin \(d_0\) and a \(P_D\)-bin, both trajectory data. Write \(W_G=\pi(d_0)D^2\) for a group \(G\). The direction width here and the error in \Cref{a06:horizontal-equation} are taken smaller than \(D\) by a power. For a parameter center \((p_c,q_c)\) and direction representative \(v_{d_0}\), positions are within \(K D\) of the moving center
\[
 z_c(t)=(t,p_c+t v_{d_0}),\qquad Y_c(t)=q_c-tp_c
\]
on occurrences of the group. Thus their integrated mass over any required time interval \(I\) of fixed power length (also scale one), in an \(x\)-bin of width \(\varrho\), \(a_s\le\varrho\le1\), costs
\[
 K |I| W_G \varrho^d .
\]
Sum the joint caps at fixed \(d_0\) on much finer base meshes to integrate over the moving spatial domain of area \(\le K D^2\) per unit time. A band of width \(K\varrho\) with \(t\)-dependent polynomial center of bounded degree and derivative \(\le K\) works too. Grids/enlargements and interpolation cost \(K\). In particular \(\mu(G)\le K W_G\), since \(x\) is bounded by \(K\). Conversely one can discard groups of mass too small compared with \(W_G/K\), and groups of too small mass fraction relative to \(\nu^\flat(G)\), at negligible cost by \(\sum_G W_G\le K\) over occurring groups (\(P\) bounded by \(K\)), and disjointness. On remaining groups active mass and prior mass therefore compare up to \(K\) with \(W_G\).

\paragraph{Counting support by whole true boxes.}
Before the group restrictions, discard true \(Q_s\)'s of current
mass less than \(m_{Q_s}/K_1\), and freeze the retained witnesses.
The budget \Cref{a06:horizontal-budgets} makes the loss at most
\(1/K_1\).  These are witnesses from the bounded-residual law
after choosing \(L\), so its pure base cap applies to them.
The full assignment supplies the reference budget, not the witness
measure.  Later contact with a group will locate the entire box
containing these witnesses.

Fix a time bin of length \(a_s\).  Each contributing true box
meets the group's reference motion to \(KD\).  In a true time
block of length \(H_s\), all these boxes lie in one enlarged
horizontal box.  To verify this, note that the reference motion
satisfies \(Y_c'=Jz_c\cdot(1,v_{d_0})\) exactly and that
\(D\le h_s^2\).  At contact the error therefore fits both
horizontal and transverse widths.  Use the plane with row
\(Jz_c(t_*)\) at the reference block midpoint.  Each meeting
box has horizontal extent \(Kh_s\); changing its centered plane
to this reference plane costs at most \(Kh_s^2\).  Thus its
whole domain lies in the common enlargement, whose base volume is
at most \(Kh_s^4\).

The pure base cap on the frozen bounded-residual law is \(K\)
per unit volume, by its unit scalar cap and \(|x|\le K\).
Each contributing true label has disjoint whole-domain witnesses
of mass at least \(h_s^4a_s^d/K\).  Their number in a meeting
time block is consequently at most \(Ka_s^{-d}\).
Only boundedly many true time blocks meet the \(a_s\)-bin,
up to rounding losses.  Each label has \(K\) affine entries.
Its ordinary slope bounds, \(D\le a_s\), and the bounded
speed of the reference motion imply that each entry supplies only
\(K\) scalar bins on the group's occurrences during this time
bin.  Over all time bins, the \((t,x)\)-support therefore needs
at most \(Ka_s^{-1-d}\) squares of side \(a_s\).

\paragraph{From support to heavy coefficient bins.}
Normalize one typical group's weights by \(W_G\).  The band cap
bounds incidence mass with affine \(x\)-coefficients in a ball
of radius \(\varrho\ge a_s\) by \(K\varrho^d\); for large
\(\varrho\), use the total mass bound.  On any dense residual,
discard trajectories carrying too little integrated incidence mass
relative to their prior.  The total prior is at most \(K\), so
a sufficiently small reciprocal-subpower threshold loses negligible
mass.  On the remaining trajectories the prior itself satisfies the
coefficient-ball cap \(K\varrho^d\).

Apply \Cref{lem:scalar-clustering} with \(p=a_s\), zero quadratic
coefficient, and the support count just proved.  It finds a prior
mass of at least \(a_s^d/K\) in an ordinary coefficient bin,
after splitting a possible subpower enlargement.  The duration
threshold gives the same lower bound, up to \(K\), for the
residual active mass in that bin.

This conclusion holds on every dense residual, so all but negligible
mass can be assigned to heavy bins.  Explicitly, if for a fixed
\(\delta>0\) the bins of mass below
\(W_Ga_s^dN^{-\delta}\) carried a nonvanishing fraction of the
dense total success, summation would select one group where this
bad-bin residual is dense relative to \(W_G\).  Applying the
preceding clustering argument only to those bins is a contradiction.
Let \(\delta\) decrease slowly.  We may therefore discard whole
bins of mass below \(W_Ga_s^d/K\), retain dense mass, and bound
the number of remaining coefficient bins per group by \(Ka_s^{-d}\).
The assignments are full-time coefficient assignments within the
group, so they predict \(x\) to \(Ka_s\) throughout unit time.

Finally subdivide time into bins of fixed-power length
\(\eta\ll a_s\).  Delete coefficient entries of mass less than
\(\eta W_Ga_s^d/K'\) in a subbin.  The number of coefficient
bins is at most \(Ka_s^{-d}\), so summing these thresholds over
the time partition and the groups gives negligible loss for
sufficiently large subpower \(K'\).  Store the retained subbin
masses for the following count.

Fix a group, a time subbin, and an \(x_{a_s}\)-bin.  If a
trajectory \(\ell_0\) in a coefficient entry meets that scalar
bin at \(t_0\), every trajectory \(\ell\) in the same
coefficient entry and every time \(t\) in the subbin satisfy
\[
 |x_\ell(t)-x_{\ell_0}(t_0)|
 \le Ca_s+K|t-t_0|\le Ca_s+K\eta\le Ka_s.
\]
Thus the entry's entire stored subbin mass lies in the enlarged
scalar band.  The time/band cap is \(K\eta W_Ga_s^d\);
summing the subbin floors beneath it gives at most \(K\)
compatible coefficient entries.

\end{proof}
